\chapter{Instructions vectorielles et code parallèle sur les données}\label{ch:ll:vector-instructions-and-data-parallel-code}

Un analyseur qui lit un octet à la fois regarde un message comme on lit une page lettre par lettre. En 2019, Geoff
Langdale et Daniel Lemire ont publié un analyseur JSON qui examine 32 ou 64 octets par instruction, et l'ont présenté comme
le premier analyseur JSON conforme à la norme «~capable de traiter des gigaoctets de données par seconde sur un seul cœur~», en utilisant un quart
ou moins des instructions d'un analyseur de référence très répandu. Sur l'ordinateur portable de ce livre, la boucle qui cherche la fin
d'un champ octet par octet parcourt environ quatre octets par nanoseconde~; la même recherche écrite avec des \omterm{def:ll:vector-instructions-and-data-parallel-code:simd}{instructions vectorielles} de 256 bits parcourt
environ quatre-vingts octets, soit une dizaine de fois plus vite, et découpe un ordre FIX de 159 octets en ses seize champs en environ \qty{11}{\nano\second}
contre une centaine. Les données de marché et les messages d'ordres sont du texte ou du binaire compacté, lus champ par champ à chaque tick~: ce
chapitre montre ce que sont les \omterm{def:ll:vector-instructions-and-data-parallel-code:simd}{instructions vectorielles}, quand le compilateur les utilise de lui-même et ce qui l'en empêche, comment les écrire à la main
avec un repli scalaire auquel les tests les confrontent, et comment disposer les données pour qu'elles s'appliquent.

\section{Registres SIMD et jeux d'instructions}

\begin{definition}[SIMD, instruction vectorielle]\label{def:ll:vector-instructions-and-data-parallel-code:simd}
Le \emph{SIMD}\index{SIMD} (single instruction, multiple data) est l'exécution d'une même opération sur plusieurs éléments de données
à la fois. Une \emph{instruction vectorielle}\index{instruction vectorielle} est une telle opération appliquée à un registre vectoriel qui contient plusieurs
éléments côte à côte~: sur x86-64, des registres \texttt{xmm} de 16 octets (SSE2, présent sur tout processeur x86-64), des
registres \texttt{ymm} de 32 octets (AVX2) et, sur les processeurs qui le proposent, des registres \texttt{zmm} de 64 octets (AVX-512).
\end{definition}

Un registre de 32 octets contient 32 octets, 16 entiers de deux octets, 8 entiers ou flottants de quatre octets, ou 4 flottants de huit octets, et une seule instruction
additionne, compare ou multiplie tous ces éléments. L'instruction la plus importante pour l'analyse syntaxique est la comparaison d'octets~:
\texttt{VPCMPEQB} compare les 32 octets de deux registres \texttt{ymm} et met chaque octet du résultat à tous ses bits à un là où ils
sont égaux, et \texttt{VPMOVMSKB} rassemble le bit de poids fort de chacun de ces octets dans un entier de 32 bits placé dans un registre ordinaire. Un
fragment de message comparé à un registre entièrement rempli de l'octet délimiteur devient donc un masque de 32 bits, avec un bit par position
et «~où se trouve le premier délimiteur~» devient «~combien de zéros de poids faible le masque comporte-t-il~», soit une instruction de plus
(\cref{fig:ll:vector-instructions-and-data-parallel-code:mask}). La boucle fait en trois instructions ce que la boucle scalaire
fait en 32 comparaisons et 32 branchements.

\begin{omfigure}
\begin{tikzpicture}[font=\scriptsize,c/.style={draw=gray!60,minimum width=0.34cm,minimum height=0.42cm,inner sep=0pt}]
\foreach \ch [count=\j from 0] in {3,5,=,D,|,5,5,=,E,S,Z,6,|,5,4,=,1,|,3,8,=,5,|,4,4,=,5,7,2,3,|,6} {
  \pgfmathtruncatemacro{\i}{\j}
  \pgfmathtruncatemacro{\hit}{(\i==4)||(\i==12)||(\i==17)||(\i==22)||(\i==30)}
  \ifnum\hit=1
    \node[c,fill=omThm!25] at (0.36*\i,0) {\ch};
    \node[c,fill=omThm!55,text=white] at (0.36*\i,-1.3) {1};
  \else
    \node[c,fill=gray!8] at (0.36*\i,0) {\ch};
    \node[c] at (0.36*\i,-1.3) {0};
  \fi
  \node[c,fill=omDef!12] at (0.36*\i,-0.65) {|};
}
\node[anchor=east,font=\footnotesize] at (-0.3,0) {32 octets d'un message};
\node[anchor=east,font=\footnotesize] at (-0.3,-0.65) {délimiteur, diffusé};
\node[anchor=east,font=\footnotesize] at (-0.3,-1.3) {bits de \texttt{movemask}};
\foreach \i in {0,8,16,24,31} {\node[font=\tiny,gray] at (0.36*\i,0.38) {\i};}
\draw[-Stealth,thick,omThm] (0.36*4,-1.55) -- ++(0,-0.45) node[below,align=center,font=\footnotesize]
  {zéros de poids faible $= 4$~:\\ premier délimiteur};
\node[anchor=west,align=left,font=\footnotesize] at (3.4,-2.25) {effacer le bit de poids faible et recompter~: 12, 17, 22, 30\\
  (un \texttt{VPCMPEQB}, un \texttt{VPMOVMSKB}, un comptage par champ)};
\end{tikzpicture}

\omcaption{Recherche des délimiteurs de champs d'un message FIX par une seule comparaison de 32 octets. L'octet \texttt{|} représente le
délimiteur SOH (octet 1). La comparaison met un octet à tous les bits à un là où le message est égal au délimiteur~; l'instruction de masque
rassemble un bit par octet~; le comptage des zéros de poids faible donne chaque position
à son tour.}\label{fig:ll:vector-instructions-and-data-parallel-code:mask}
\end{omfigure}

Le processeur de l'ordinateur portable dispose de SSE2 jusqu'à AVX2 et de la multiplication-addition fusionnée (FMA) (livre~4, chapitre~25), mais pas d'AVX-512, comme le montrent ses
indicateurs de fonctionnalités dans \path{/proc/cpuinfo}. C'est la situation ordinaire pour le code d'une société de trading~: il est compilé une fois et
s'exécute sur des machines de plusieurs générations, si bien que le jeu d'instructions le plus large ne peut pas être présumé, et un binaire qui exécute une
instruction AVX2 sur un processeur qui ne la possède pas s'arrête avec un signal d'instruction illégale. La section~3 traite ce point.

\section{Auto-vectorisation et ce qui la bloque}

\begin{definition}[Auto-vectorisation]\label{def:ll:vector-instructions-and-data-parallel-code:auto}
L'\emph{auto-vectorisation}\index{auto-vectorisation} est la transformation, par le compilateur, d'une boucle scalaire en \omterm{def:ll:vector-instructions-and-data-parallel-code:simd}{instructions vectorielles}
qui traitent plusieurs itérations à la fois, effectuée lorsqu'il peut démontrer que le résultat reste inchangé et juge la
transformation rentable.
\end{definition}

Le compilateur vectorise de lui-même plus souvent qu'on ne le pense, et moins souvent qu'on ne l'espère. GCC le fait avec
\texttt{-O3}~; à partir de GCC~12, il le fait aussi à \texttt{-O2}, avec un modèle de coût qui n'accepte que des transformations très bon marché
(les notes de version), si bien qu'une boucle que GCC~11 laisse scalaire à \texttt{-O2} peut être vectorisée par GCC~13 avec le
même drapeau~: la version du compilateur fait partie de la performance. \texttt{-fopt-info-vec-all} demande au compilateur de signaler chaque décision.
Le \Cref{lst:ll:vector-instructions-and-data-parallel-code:report} est son rapport sur cinq petites boucles avec le GCC~11 de l'ordinateur portable.

\omcode{code/low-latency/14-vector-instructions-and-data-parallel-code/cpp/vecreport/gxx11.txt}{1}{20}{Le rapport de vectorisation
de GCC~11 pour \texttt{ll\_vecdemo.cpp}~: la réévaluation sur des colonnes (ligne 13), sur des enregistrements (20--21), une recherche qui
s'arrête tôt (27), et des sommes entière et en double précision (34, 40).}\label{lst:ll:vector-instructions-and-data-parallel-code:report}

Chaque ligne est une leçon. À \texttt{-O2}, rien n'est vectorisé. La réévaluation sur des tableaux séparés est vectorisée à
\texttt{-O3}, deux flottants en double précision par instruction, quatre avec \texttt{-mavx2}, mais seulement après que le compilateur a «~versionné~» la boucle~:
le pointeur de sortie pourrait chevaucher une entrée, il le vérifie donc à l'exécution et conserve une copie scalaire de la boucle pour ce cas
(\texttt{\_\_restrict}, qui promet l'absence de chevauchement, supprime la vérification). Le même calcul sur des enregistrements de 80 octets n'est
pas vectorisé du tout~: les cinq champs d'options consécutives sont séparés de 80 octets, et GCC~11 ne trouve aucun type vectoriel pour de tels
chargements à pas variable. La recherche est refusée parce qu'elle peut quitter la boucle plus tôt («~control flow in loop~»)~: GCC~11 ne vectorise pas une boucle
avec sortie anticipée. La somme entière est vectorisée. La somme en double précision est signalée
vectorisée elle aussi, et son assembleur montre pourquoi cela ne signifie rien~: les additions restent dans leur ordre d'origine, un élément à la
fois, car l'addition en virgule flottante n'est pas associative et en modifier l'ordre changerait le résultat (chapitre~8, où
\texttt{-ffast-math} l'autorisait et modifiait la réponse). Les obstacles, en bref~: le recouvrement possible des pointeurs, les données à pas variable ou dispersées,
les sorties anticipées, les appels qui ne peuvent pas être mis en ligne, et les réductions dont l'ordre compte.

\section{Intrinsèques et répartition selon les fonctions du CPU}

\begin{definition}[Intrinsèque du compilateur, répartition selon les fonctions du CPU]\label{def:ll:vector-instructions-and-data-parallel-code:intrinsic}
Un \emph{intrinsèque du compilateur}\index{intrinsèque du compilateur} est une fonction que le compilateur connaît et traduit en une instruction précise
(\texttt{\_mm256\_cmpeq\_epi8} en \texttt{VPCMPEQB}), ce qui permet à un programme d'utiliser une instruction sans écrire
d'assembleur, tandis que le compilateur continue d'allouer les registres et d'ordonnancer autour d'elle. La \emph{répartition selon les
fonctions du CPU}\index{répartition selon les fonctions du CPU} choisit à l'exécution, d'après les fonctions que signale le processeur, laquelle parmi plusieurs versions compilées d'une fonction
doit être appelée.
\end{definition}

Lorsque le compilateur refuse de vectoriser une boucle, les intrinsèques le font à la main. La recherche d'un octet dans la construction charge 32 octets,
les compare au délimiteur diffusé, prend le masque et compte ses zéros de poids faible~; un message dont la longueur n'est pas un multiple de
32 se termine par un chargement de plus qui s'achève exactement sur son dernier octet et chevauche le bloc précédent, les octets déjà
vus étant écartés du masque, de sorte qu'aucun octet au-delà du message n'est jamais lu. Ce dernier point n'est pas une affaire de pédanterie~: une lecture au-delà
de la fin d'un tampon qui se termine à une frontière de page provoque une faute, et le sanitiseur du chapitre~8 signale toute lecture de ce type.

\omcode{code/firm/simdscan/cpp/firm_simdscan.hpp}{64}{79}{Recherche AVX2 d'un octet, avec un dernier
bloc qui chevauche le précédent.}\label{lst:ll:vector-instructions-and-data-parallel-code:find}

\begin{sloppypar}La fonction est compilée pour AVX2 au moyen d'un attribut \texttt{target}, sans compiler pour autant tout le programme, et elle est appelée
seulement après une vérification. La fonction \texttt{\_\_builtin\_cpu\_supports(\textquotedbl{}avx2\textquotedbl{})} de GCC renvoie une valeur non nulle lorsque le processeur en fonctionnement
dispose de la fonctionnalité~; la construction la vérifie une seule fois, au premier appel, et en conserve la réponse. Le \omterm{def:ll:cpp-for-latency-iii-the-compiler:fmv}{multiversionnage de fonctions} du chapitre~8 fait la même chose
automatiquement pour une fonction entière~; la version explicite rend le choix visible et testable, et les tests appellent chaque
version directement sur les mêmes entrées.\end{sloppypar}

\omcode{code/firm/simdscan/cpp/firm_simdscan.hpp}{196}{210}{Répartition~: détecter les fonctions une seule fois, puis appeler la
version la plus large.}\label{lst:ll:vector-instructions-and-data-parallel-code:dispatch}

Rust exprime le même contrat dans ses types. Une fonction qui utilise des intrinsèques est déclarée \texttt{unsafe} avec l'attribut \texttt{target\_feature} qui active \texttt{avx2}, et son appel est sûr
seulement une fois que la macro \texttt{is\_x86\_feature\_detected!} a renvoyé vrai pour \texttt{avx2}, comme le montre l'exemple même de la bibliothèque standard~: le
bloc \texttt{unsafe} «~est sûr parce que nous testons que la fonctionnalité \texttt{avx2} est bien disponible sur notre CPU~».

\omcode{code/firm/simdscan/rust/src/lib.rs}{86}{98}{La même répartition en Rust~: l'appel unsafe est protégé par la
vérification à l'exécution.}\label{lst:ll:vector-instructions-and-data-parallel-code:rust}

Les nombres demandent un second artifice. Huit chiffres ASCII chargés comme un seul entier de 64 bits peuvent être convertis en leur valeur sans boucle~:
on soustrait le caractère \texttt{0} de chaque octet en une seule opération, puis on combine les chiffres voisins en nombres à deux chiffres dans des voies de 16 bits
ceux-ci en nombres à quatre chiffres dans des voies de 32 bits, et ceux-là dans le résultat, soit trois multiplications en tout. Il s'agit de \omterm{def:ll:vector-instructions-and-data-parallel-code:simd}{SIMD}
à l'intérieur d'un registre ordinaire, et cela ne demande aucune instruction particulière~; une version avec les multiplications-additions de SSE4.1 fait la même chose dans les
registres vectoriels.

\omcode{code/firm/simdscan/cpp/firm_simdscan.hpp}{134}{144}{Huit chiffres en trois étapes de multiplication-addition dans un seul
registre de 64 bits.}\label{lst:ll:vector-instructions-and-data-parallel-code:swar}

\begin{proposition}[La combinaison des huit chiffres]\label{prop:ll:vector-instructions-and-data-parallel-code:swar}
Soient des octets $b_0,\dots,b_7$ contenant les chiffres $d_0,\dots,d_7$ (après soustraction du caractère zéro), $b_0$ étant le plus bas. Après
l'étape $v \leftarrow (10v + (v \gg 8)) \wedge \mathtt{0x00FF00FF00FF00FF}$, la voie de 16 bits $k$ contient $10 d_{2k} + d_{2k+1}$~;
après les deux étapes analogues avec 100 et 10\,000, les 32 bits de poids faible contiennent $\sum_{j=0}^{7} d_j\, 10^{7-j}$.
\end{proposition}

\begin{proof}
Avant le masque, l'octet $2k$ de $10v + (v \gg 8)$ vaut $10 d_{2k} + d_{2k+1} \le 99 < 256$, puisque le décalage ramène l'octet $2k+1$
en position $2k$~; aucune retenue ne sort de l'octet, et le masque écarte les octets impairs, qui contiennent des sommes inutiles. Le même
raisonnement sur les voies de 16 bits (valeurs allant jusqu'à $99 \times 100 + 99 < 2^{16}$) et sur les voies de 32 bits (jusqu'à $9\,999 \times 10\,000 +
9\,999 < 2^{32}$) donne les valeurs à quatre et à huit chiffres, le chiffre antérieur occupant toujours la position la plus élevée.
\end{proof}

Sur le portable, un champ de huit chiffres coûte environ \qty{2.0}{\nano\second} avec la boucle scalaire, \qty{1.0}{\nano\second} avec le
procédé du registre et \qty{0.7}{\nano\second} avec SSE4.1, contre environ \qty{5}{\nano\second} pour la
bibliothèque standard, \texttt{std::from\_chars}, qui valide en outre et traite une longueur quelconque. Un prix converti en un nombre entier de
dix-millièmes coûte environ \qty{7}{\nano\second}, et le même prix lu comme un double par \texttt{strtod} environ
\qty{43}{\nano\second}. L'article de Lemire sur l'analyse des nombres illustre ce point pour l'analyse des flottants~: un flottant de 64 bits comporte au plus
17 chiffres significatifs, qui tiennent dans un seul mot de 64 bits, et une analyse faite ainsi est plusieurs fois plus rapide que celle de la bibliothèque C.

\section{Où la vectorisation est rentable dans un système de trading}

La \cref{fig:ll:vector-instructions-and-data-parallel-code:scan} mesure la recherche d'un délimiteur placé à la
fin d'un message de chaque longueur. La boucle scalaire coûte environ un quart de nanoseconde par octet~; SSE2 environ un quarantième, AVX2 environ un
quatre-vingtième, auxquels s'ajoute une fraction de nanoseconde par appel. En dessous de 16 octets, les deux versions vectorielles confient le message à la boucle scalaire
et coûtent la même chose. Découper un message d'ordre FIX complet en ses champs (16 champs, 159 octets) prend environ
\qty{95}{\nano\second} un octet à la fois et environ \qty{11}{\nano\second} avec la boucle à masque, soit environ neuf fois moins.

\begin{omfigure}
\begin{tikzpicture}
\begin{axis}[width=11.5cm,height=5.8cm,xmode=log,ymode=log,log basis x=2,xmin=6,xmax=5500,ymin=0.4,ymax=2000,
  xlabel={longueur du message (octets, échelle logarithmique)},ylabel={ns par recherche (échelle logarithmique)},
  xtick={8,16,32,64,128,256,512,1024,4096},xticklabels={8,16,32,64,128,256,512,1024,4096},
  tick label style={font=\small},label style={font=\small},grid=major,grid style={gray!25},
  legend image code/.code={\draw[#1] (0cm,0cm) -- (0.6cm,0cm);},
  legend columns=3,legend style={font=\footnotesize,at={(0.5,-0.3)},anchor=north,draw=none,fill=none,
  /tikz/every even column/.append style={column sep=8pt}},table/col sep=comma]
\addplot[omThm,very thick,mark=*,mark size=1.2pt] table[x=bytes,y=scalar]{figdata/low-latency/14-vector-instructions-and-data-parallel-code/measured_scan.csv};
\addplot[omMeth,thick,mark=square*,mark size=1.2pt] table[x=bytes,y=sse2]{figdata/low-latency/14-vector-instructions-and-data-parallel-code/measured_scan.csv};
\addplot[omDef,very thick,mark=triangle*,mark size=1.6pt] table[x=bytes,y=avx2]{figdata/low-latency/14-vector-instructions-and-data-parallel-code/measured_scan.csv};
\legend{boucle scalaire,SSE2 (16 octets),AVX2 (32 octets)}
\end{axis}
\end{tikzpicture}

\omcaption{Temps nécessaire pour trouver un délimiteur à la fin d'un message, selon la longueur du message~: médiane de lots de mille recherches
d'un message dans le cache de premier niveau. Mesuré sur un portable (Intel Core Ultra 7 155H) sous WSL2, sans cœurs isolés. Données~:
\texttt{bench\_simd.py}.}\label{fig:ll:vector-instructions-and-data-parallel-code:scan}
\end{omfigure}

Où cela compte-t-il~? Partout où des octets sont examinés un à un sur le chemin critique~: découper les messages FIX sur leur délimiteur
(chapitre~15), repérer les caractères structurels des messages JSON provenant des plateformes qui les publient dans ce format (chapitre~17), valider et
analyser les champs numériques, les sommes de contrôle et les copies. Cela compte moins pour les protocoles binaires dont les champs sont à des décalages fixes
(chapitre~16), où il n'y a rien à rechercher~; là, les \omterm{def:ll:vector-instructions-and-data-parallel-code:simd}{instructions vectorielles} rapportent plutôt dans le traitement en masse de tableaux~: les
contrôles de risque de nombreux ordres, la réévaluation d'un livre de positions, les statistiques qu'une stratégie met à jour. Trois conditions décident~:
assez d'éléments pour remplir un registre, la même opération sur chacun, et des données disposées de façon contiguë. Les deux premières sont des propriétés
du problème~; la troisième relève d'un choix de conception.

\section{Structure de tableaux}

\begin{definition}[Structure de tableaux]\label{def:ll:vector-instructions-and-data-parallel-code:soa}
Une \emph{structure de tableaux}\index{structure de tableaux} stocke une collection d'enregistrements sous la forme d'un tableau par champ, de sorte que le
même champ d'enregistrements consécutifs soit contigu en mémoire~; la disposition habituelle, un enregistrement après l'autre, s'appelle un
tableau de structures.
\end{definition}

Un livre d'options réévalué à partir de ses grecques montre la différence
(\cref{fig:ll:vector-instructions-and-data-parallel-code:layout}). Chaque enregistrement d'option porte un identifiant, un symbole, sa
quantité et quatre grecques, une échéance et quelques indicateurs~: 80 octets, dont la réévaluation $q(\Delta\,\delta S + \tfrac12
\Gamma\,\delta S^2 + \mathcal{V}\,\delta\sigma + \Theta\,\delta t)$ n'en lit que 40. Stockés en enregistrements, quatre deltas consécutifs sont
séparés de 80 octets et aucune instruction ne les charge ensemble~; stockés en colonnes, ils occupent 32 octets consécutifs, en un seul chargement. La moitié de
chaque \omterm{def:ll:memory-hierarchy-and-caches:line}{ligne de cache} récupérée pour les enregistrements est de surcroît perdue.

\begin{omfigure}
\begin{tikzpicture}[font=\scriptsize]
\node[anchor=west,font=\footnotesize\bfseries] at (-0.1,1.0) {tableau de structures (enregistrements de 80 octets)};
\foreach \r in {0,1,2} {
  \pgfmathsetmacro{\x}{\r*4.3}
  \draw[fill=gray!10] (\x,0) rectangle ++(0.8,0.5) node[midway]{id};
  \draw[fill=gray!10] (\x+0.8,0) rectangle ++(1.1,0.5) node[midway,font=\tiny]{symbole};
  \foreach \f/\k in {q/0,$\Delta$/1,$\Gamma$/2,$\mathcal{V}$/3,$\Theta$/4} {
    \draw[fill=omDef!25] (\x+1.9+0.4*\k,0) rectangle ++(0.4,0.5) node[midway]{\f};
  }
  \draw[fill=gray!10] (\x+3.9,0) rectangle ++(0.35,0.5);
}
\foreach \r in {0,1,2} {\draw[-Stealth,omThm,thick] (\r*4.3+2.5,-0.05) -- (\r*4.3+2.5,-0.45);}
\node[anchor=west,omThm] at (-0.1,-0.7) {deltas espacés de 80 octets~: un chargement chacun, la moitié de chaque ligne inutilisée};
\node[anchor=west,font=\footnotesize\bfseries] at (-0.1,-1.45) {structure de tableaux (un tableau par champ)};
\foreach \f/\y in {$\Delta$/-2.15,$\Gamma$/-2.75} {
  \node[anchor=east] at (-0.15,\y+0.25) {\f};
  \foreach \k in {0,...,15} {\draw[fill=omDef!25] (0.4*\k,\y) rectangle ++(0.4,0.5);}
}
\draw[very thick,omThm] (0,-2.15) rectangle ++(1.6,0.5);
\node[anchor=west,omThm] at (6.6,-1.9) {quatre deltas, 32 octets contigus~:};
\node[anchor=west,omThm] at (6.6,-2.3) {un chargement \texttt{ymm}, une multiplication-addition fusionnée (FMA)};
\end{tikzpicture}

\omcaption{Le même carnet d'options, en enregistrements puis en colonnes. La réévaluation lit cinq champs de l'enregistrement~; en colonnes, chaque
champ de quatre options consécutives remplit un registre de 256~bits.}\label{fig:ll:vector-instructions-and-data-parallel-code:layout}
\end{omfigure}

La \Cref{fig:ll:vector-instructions-and-data-parallel-code:reval} mesure les deux dispositions pour trois jeux d'indicateurs. Avec 1\,024
options, qui tiennent dans les caches, les enregistrements coûtent environ \qty{0.8}{\nano\second} par option, quels que soient les indicateurs, car aucun
compilateur ne les vectorise~; les colonnes coûtent le même prix avec \texttt{-O2} (GCC~11 ne les vectorise pas à ce niveau), la moitié avec \texttt{-O3}
(deux doubles par instruction) et le quart avec \texttt{-mavx2 -mfma}, quatre doubles par instruction et une multiplication-addition
fusionnée (FMA) pour chaque couple multiplication et addition. Avec un million d'options, les données ne tiennent plus dans les caches et les
octets déplacés décident~: 88~octets par option pour les enregistrements (l'enregistrement et le résultat), 48 pour les colonnes, et le coût mesuré est
environ trois fois plus faible pour les colonnes, quels que soient les indicateurs, la largeur des vecteurs ne jouant plus guère.

\begin{omfigure}
\begin{tikzpicture}
\begin{axis}[width=12cm,height=5.8cm,ybar,bar width=8pt,ymin=0,ymax=6.5,
  symbolic x coords={O2,O3,O3avx2},xtick=data,xticklabels={\texttt{-O2},\texttt{-O3},\texttt{-O3 -mavx2 -mfma}},
  enlarge x limits=0.25,ylabel={ns par option},tick label style={font=\small},label style={font=\small},grid=major,
  grid style={gray!25},legend columns=2,legend style={font=\footnotesize,at={(0.5,-0.22)},anchor=north,draw=none,fill=none,
  /tikz/every even column/.append style={column sep=8pt}},table/col sep=comma]
\addplot[fill=gray!35,draw=gray] table[x=key,y=rows_1k]{figdata/low-latency/14-vector-instructions-and-data-parallel-code/measured_reval.csv};
\addplot[fill=omDef!70,draw=omDef] table[x=key,y=cols_1k]{figdata/low-latency/14-vector-instructions-and-data-parallel-code/measured_reval.csv};
\addplot[fill=gray!70,draw=gray!80!black] table[x=key,y=rows_1m]{figdata/low-latency/14-vector-instructions-and-data-parallel-code/measured_reval.csv};
\addplot[fill=omDef!35,draw=omDef] table[x=key,y=cols_1m]{figdata/low-latency/14-vector-instructions-and-data-parallel-code/measured_reval.csv};
\legend{{enregistrements, 1\,024 options},{colonnes, 1\,024 options},{enregistrements, $2^{20}$ options},{colonnes, $2^{20}$ options}}
\end{axis}
\end{tikzpicture}

\omcaption{Réévaluation d'un carnet d'options à partir de ses grecques, stockées en enregistrements de 80~octets ou en colonnes, compilées avec trois jeux
d'indicateurs par GCC~11~: en cache (1\,024 options) et depuis la mémoire ($2^{20}$ options, 80~Mo d'enregistrements). Mesures sur un portable (Intel Core
Ultra 7 155H) sous WSL2, sans cœurs isolés. Données~: \texttt{bench\_simd.py}.}\label{fig:ll:vector-instructions-and-data-parallel-code:reval}
\end{omfigure}

Les colonnes ont un coût. Ajouter ou retirer une option touche chaque tableau~; le code qui traite une option à la fois (un ordre
qui arrive, une exécution) lit un élément dans chacun de plusieurs tableaux, ce qui fait plus de \omterm{def:ll:memory-hierarchy-and-caches:line}{lignes de cache} qu'un seul enregistrement~; et la
disposition est moins naturelle à écrire. Le compromis habituel consiste à garder les enregistrements là où l'on traite des éléments isolés, et les colonnes
là où l'on calcule sur de nombreux éléments, ou des blocs de colonnes (les deltas de huit options, puis leurs gammas) qui offrent les deux. Un carnet d'ordres
traité un message à la fois (chapitre~19) demande des enregistrements~; un moteur de risque qui réévalue chaque position à chaque mouvement (chapitre~22)
demande des colonnes.

\section{Tutoriel~: parcourir, analyser et disposer}\label{tut:ll:vector-instructions-and-data-parallel-code:tut}

\begin{tutorial}
\textbf{Objectif.} Parcourir et analyser avec des \omterm{def:ll:vector-instructions-and-data-parallel-code:simd}{instructions vectorielles} derrière un contrôle des fonctions du CPU, vérifiées octet par octet contre une référence scalaire~; lire
les décisions du compilateur~; mesurer les dispositions. \textbf{État final~:}
\cref{fig:ll:vector-instructions-and-data-parallel-code:scan,fig:ll:vector-instructions-and-data-parallel-code:reval},
\cref{lst:ll:vector-instructions-and-data-parallel-code:report}, et des tests verts en C++ et en Rust.
\begin{tutsteps}
\item \textbf{Vérifier le code vectoriel par rapport au code scalaire.} Le test C++ de la construction parcourt toutes les longueurs de 0 à 200 octets avec
des délimiteurs placés à des positions aléatoires (ou aucun), vingt messages par longueur, à travers les recherches scalaire, SSE2, AVX2 et réparties
ainsi que les recherches de positions, et exige des réponses identiques~; il analyse 200\,000 nombres à huit chiffres de quatre façons. Les
tampons ont exactement la longueur des messages, si bien qu'une exécution sous AddressSanitizer détecte toute lecture au-delà de la fin.
\item \textbf{Lire le rapport du compilateur}~: \texttt{python ll\_vecreport.py} compile \texttt{ll\_vecdemo.cpp} avec trois jeux
d'indicateurs et écrit le \cref{lst:ll:vector-instructions-and-data-parallel-code:report} pour le GCC installé~; l'exécuter avec un
GCC plus récent et comparer les lignes de \texttt{-O2}.
\item \textbf{Mesurer} avec \texttt{python bench\_simd.py}~: recherches selon la longueur, répartition des messages FIX, analyseurs et
réévaluation construite avec trois jeux d'indicateurs.
\item \textbf{La même chose en Rust}~: \texttt{cargo test} dans \path{code/firm/simdscan/rust} exécute les parcours de la réplique pour chaque longueur
et ses analyseurs sur les mêmes valeurs.
\end{tutsteps}
\textbf{Pour aller plus loin.} Ajouter \texttt{\_\_restrict} aux pointeurs de \texttt{reval\_columns} et vérifier que la
ligne de versionnement du rapport disparaît~; écrire la somme des doubles avec quatre accumulateurs distincts et voir si le compilateur
la vectorise pour de bon.
\end{tutorial}

\section{Construction~: le scanneur de la société}\label{bld:ll:vector-instructions-and-data-parallel-code:simdscan}

\begin{build}
\bfield{Rôle} Les primitives octet par octet de chaque protocole texte que la société lit~: le moteur FIX découpe les messages sur le
délimiteur avec elles (chapitre~15), et le client WebSocket repère les caractères structurels du JSON et analyse les prix et les tailles avec
elles (chapitre~17).
\bfield{Interface} C++20 \texttt{firm::simdscan}~: \texttt{find\_byte(p, n, c)}, \texttt{positions(p, n, c, out, cap)},
\texttt{positions\_any(p, n, set, m, out, cap)}, chacune en version \texttt{\_scalar}, \texttt{\_sse2} (recherche) et \texttt{\_avx2}
~; \texttt{parse8\_scalar}, \texttt{parse8\_swar}, \texttt{parse8\_sse}, \texttt{is8digits},
\texttt{parse\_uint(p, n, value, used)}, \texttt{parse\_fixed(p, n, scale, out)}~; \texttt{level()}. Rust
\texttt{firm\_simdscan}~: \texttt{find\_byte}, \texttt{positions}, \texttt{parse8\_swar}, \texttt{parse\_uint},
\texttt{parse\_fixed}, avec les références scalaires.
\bfield{Règles} Chaque routine vectorielle renvoie ce que sa référence scalaire renvoie pour toute entrée, ne lit jamais en dehors de
$[p, p+n)$, et n'est appelée qu'après le contrôle des fonctions du CPU à l'exécution~; le code AVX2 est compilé par des attributs de fonction, jamais en
compilant le programme pour AVX2~; les prix sont analysés en entiers d'une échelle déclarée, jamais en doubles.
\bfield{Tests d'acceptation} \texttt{code/firm/simdscan/}~: chaque longueur de 0 à 200 octets contre les références scalaires en
C++ et en Rust, la limite de capacité des recherches de positions respectée, 200\,000 valeurs à huit chiffres analysées de façon identique par chaque
méthode, entiers de 19 chiffres, prix avec et sans décimales, prix négatifs, et le rejet des prix mal formés.
\bfield{Pour aller plus loin} Une version AVX-512 sélectionnée par le même répartiteur~; un classifieur des caractères structurels du JSON avec une
table de consultation (une permutation au lieu d'une comparaison par caractère)~; un validateur pour huit chiffres dans la voie vectorielle.
\end{build}

\begin{omsources} \sloppy
\item G.~Langdale et D.~Lemire, «~Parsing gigabytes of JSON per second~», \emph{The VLDB Journal} 28(6), 2019.
\item D.~Lemire, «~Number parsing at a gigabyte per second~», \emph{Software: Practice and Experience} 51(8), 2021.
\item Notes de version de GCC 12 (vectorisation à \texttt{-O2})~; manuel de GCC, fonctions intégrées x86
(\texttt{\_\_builtin\_cpu\_supports}).
\item Intel 64 and IA-32 Architectures Software Developer's Manual, instructions \texttt{PCMPEQB} et \texttt{PMOVMSKB} (telles
que transcrites dans la référence des instructions x86 sur felixcloutier.com).
\item La bibliothèque standard de Rust, \texttt{core::arch} et \texttt{is\_x86\_feature\_detected}.
\end{omsources}

\section{Exercices}

\begin{exercise}[$\star$]\label{exo:ll:vector-instructions-and-data-parallel-code:1}
Combien de doubles, d'entiers de 32~bits et d'octets un registre de 256~bits contient-il~? Combien de comparaisons de 256~bits faut-il pour
examiner chaque octet d'un paquet de 1\,500 octets~?
\end{exercise}

\begin{exercise}[$\star$]\label{exo:ll:vector-instructions-and-data-parallel-code:2}
Le masque d'un bloc de 32~octets vaut \texttt{0x00041010} (bit~0 pour le premier octet). Où se trouvent les délimiteurs dans le bloc~?
\end{exercise}

\begin{exercise}[$\star$]\label{exo:ll:vector-instructions-and-data-parallel-code:3}
Appliquer à la main la \cref{prop:ll:vector-instructions-and-data-parallel-code:swar} aux chiffres \texttt{20260925}~: donner les
voies de 16~bits après la première étape et les voies de 32~bits après la seconde.
\end{exercise}

\begin{exercise}[$\star\star$]\label{exo:ll:vector-instructions-and-data-parallel-code:4}
Un message de 70~octets est parcouru par blocs de 32~octets, le dernier bloc chevauchant le précédent. Indiquer les octets couverts par chaque chargement, combien
de bits il faut décaler hors du dernier masque, et pourquoi le résultat reste juste si le délimiteur se trouve à l'octet 60.
\end{exercise}

\begin{exercise}[$\star\star$]\label{exo:ll:vector-instructions-and-data-parallel-code:5}
Expliquer chacune des trois raisons que donne le compilateur dans le \cref{lst:ll:vector-instructions-and-data-parallel-code:report}
pour ne pas vectoriser, ou pour ne vectoriser qu'après un contrôle, et comment supprimer chacune d'elles.
\end{exercise}

\begin{exercise}[$\star\star$]\label{exo:ll:vector-instructions-and-data-parallel-code:6}
Estimer, à partir du balayage mesuré, le temps que met AVX2 pour parcourir un tampon entier de 4\,096~octets, et ce que cela
représente en octets par nanoseconde. Dire ce qui le limite une fois que le tampon ne tient plus dans le cache de premier niveau.
\end{exercise}

\begin{exercise}[$\star\star\star$]\label{exo:ll:vector-instructions-and-data-parallel-code:7}
\emph{Programmation.} Ajouter \texttt{positions\_any\_sse2} à \texttt{firm.simdscan} (blocs de 16~octets, dernier bloc chevauchant), puis
l'utiliser dans la \omterm{def:ll:vector-instructions-and-data-parallel-code:intrinsic}{répartition selon les fonctions du CPU} lorsque AVX2 est absent, et étendre le test C++ pour le comparer à la référence scalaire pour toutes les longueurs.
\end{exercise}

\begin{exercise}[$\star\star\star$]\label{exo:ll:vector-instructions-and-data-parallel-code:8}
\emph{Trouver la faille.} «~Notre machine de construction dispose d'AVX-512, donc nous compilons avec \texttt{-march=native}~; elle a réussi tous les tests dans
l'intégration continue et a bien fonctionné sur les nouveaux serveurs.~»
\end{exercise}

\section{Problème~: Trente-deux octets à la fois}

\begin{problem}[{Problème du week-end --- quand un balayage vectoriel est-il rentable~?}]\label{pb:ll:vector-instructions-and-data-parallel-code:1}
Le balayage mesuré (\texttt{measured\_scan.csv}) donne le coût de chaque recherche selon la longueur du message~; les mesures du découpage et de l'analyseur
figurent dans \texttt{measured\_split.csv} et \texttt{measured\_parse.csv}. Une passerelle d'ordres FIX reçoit des rapports d'exécution
d'environ 200~octets, comportant 20 champs, à raison de 200\,000 par seconde en pointe.

\medskip\noindent\textbf{Partie I --- Les coûts.}
\begin{enumerate}
\item À partir des mesures à 4\,096~octets, calculer le coût par octet de chaque méthode.
\item Ajuster une droite $a + bL$ aux coûts d'AVX2 pour $L \ge 32$ (\texttt{ll\_simd.fit}). Donner $a$ et $b$.
\item En ne comptant pour la boucle scalaire que son coût par octet, déterminer à partir de quelle longueur la recherche AVX2 est moins coûteuse, puis la même chose pour SSE2.
\item Expliquer pourquoi ne lui imputer aucun coût fixe la favorise.
\end{enumerate}

\medskip\noindent\textbf{Partie II --- La passerelle.}
\begin{enumerate}[resume]
\item Calculer le coût du découpage d'un rapport de 200~octets avec chaque méthode, par extrapolation à partir du découpage mesuré.
\item Calculer le temps de cœur consommé par seconde pour le découpage en pointe, pour chaque méthode.
\item Chaque rapport porte quatre prix et quatre quantités à analyser. En utilisant \texttt{parse\_fixed} plutôt que \texttt{strtod}, calculer
le gain réalisé par rapport.
\item Déterminer lequel des deux gains par rapport est le plus important~: le découpage ou l'analyse.
\end{enumerate}

\medskip\noindent\textbf{Partie III --- Les limites.}
\begin{enumerate}[resume]
\item Expliquer pourquoi la recherche vectorielle n'apporte rien à un protocole binaire dont les champs se trouvent à des décalages fixes.
\item Décrire ce qui se passe sur un processeur dépourvu d'AVX2.
\item Expliquer pourquoi le dernier bloc doit chevaucher le précédent plutôt que de lire 32~octets à partir du début du reste, qui dépassent la fin du tampon.
\item Dire ce qu'AVX-512 changerait dans le coût par octet, et ce qu'il ne changerait pas.
\end{enumerate}

\medskip\noindent\textbf{Partie IV --- Le verdict.}
\begin{enumerate}[resume]
\item Énoncer le \emph{résultat nommé}~: la longueur de message en dessous de laquelle le balayage vectoriel ne rentabilise pas sa mise en place, à partir du
coût mesuré par octet et par appel.
\item Vérifier si un champ FIX est plus court que ce seuil.
\item Indiquer la mesure à ajouter avant de se fier au seuil de rentabilité sur un serveur.
\item Expliquer pourquoi les références scalaires sont conservées dans la construction.
\item Décrire un test permettant de vérifier que la répartition retient le chemin scalaire sur une machine sans AVX2, puis sur une machine qui en dispose.
\item Indiquer où, ailleurs dans un système de trading, le même schéma de traitement par blocs de 32~octets s'appliquerait.
\item Indiquer quel changement de disposition permettrait aux contrôles de risque du chapitre~22 d'utiliser les mêmes instructions.
\item En une phrase~: dire quand les \omterm{def:ll:vector-instructions-and-data-parallel-code:simd}{instructions vectorielles} sont rentables.
\end{enumerate}
\end{problem}

\section{Questions d'entretien}

\begin{interviewq}[$\star$ \iqroles{développeur}]\label{iq:ll:vector-instructions-and-data-parallel-code:1}
Qu'est-ce que le \omterm{def:ll:vector-instructions-and-data-parallel-code:simd}{SIMD}~? Donnez un exemple de boucle qui en profite et d'une boucle qui n'en profite pas.
\end{interviewq}

\begin{interviewq}[$\star\star$ \iqroles{développeur}]\label{iq:ll:vector-instructions-and-data-parallel-code:2}
Comment chercheriez-vous la première occurrence d'un octet dans un tampon le plus rapidement possible~?
\end{interviewq}

\begin{interviewq}[$\star\star$ \iqroles{développeur}]\label{iq:ll:vector-instructions-and-data-parallel-code:3}
Qu'est-ce qui empêche un compilateur de vectoriser une boucle~? Comment découvririez-vous pourquoi il ne l'a pas fait~?
\end{interviewq}

\begin{interviewq}[$\star\star$ \iqroles{développeur}]\label{iq:ll:vector-instructions-and-data-parallel-code:4}
Votre binaire utilise AVX2 et doit fonctionner sur des machines plus anciennes. Comment le déployez-vous~?
\end{interviewq}

\begin{interviewq}[$\star\star$ \iqroles{développeur, chercheur}]\label{iq:ll:vector-instructions-and-data-parallel-code:5}
Tableau de structures ou \omterm{def:ll:vector-instructions-and-data-parallel-code:soa}{structure de tableaux} pour un livre de positions réévalué à chaque mouvement de prix~? Pourquoi~?
\end{interviewq}

\begin{interviewq}[$\star\star\star$ \iqroles{développeur}]\label{iq:ll:vector-instructions-and-data-parallel-code:6}
Analysez un prix décimal ASCII en un entier à virgule fixe sans boucle sur les caractères. Que devez-vous valider~?
\end{interviewq}
